%% ==================================================
%% Section: Local channel circuits and conversions
%% Sources: arXiv:2307.01696 (Malz, Styliaris, Wei, Cirac), proof of
%% Theorem 1 in the Supplemental Material; arXiv:2103.13367 (Piroli,
%% Styliaris, Cirac), paragraphs "Quantum circuits and LOCC" and "State
%% transformations with QC and LOCC".
%% ==================================================

\section{Local channel circuits and conversions}
\label{sec:qc_channels}

Lemmas~\ref{lem:qc_light_cone} and~\ref{lem:qc_light_cone_correlations}
extend from unitary gates to channels, the operations
of a local circuit that may add product ancillas and discard them
\cite{Piroli2021LOCC}.

\begin{definition}[Local channel circuit of depth $T$]
    \label{def:qc_local_channel_circuit}
    \uses{def:qc_local_circuit}
    \lean{QuantumCircuit.ChannelLayer, QuantumCircuit.BondChannelLayer,
        QuantumCircuit.BondChannelLayer.map,
        QuantumCircuit.BondChannelLayer.dual, QuantumCircuit.channelCircuitMap,
        QuantumCircuit.channelCircuitDual,
        QuantumCircuit.IsChannelPreparedInDepth,
        QuantumCircuit.IsBondChannelPreparedInDepth}
    \leanok
    Fix a bond geometry, for instance the pairs $\{k,k+1\}$ of the ring. A
    \emph{channel on the bond} $B$ is a map
    $\rho\mapsto\sum_jK_j\rho K_j^\dagger$ on density matrices of all the
    sites, where each $K_j$ is an operator on all the sites acting trivially
    off $B$ and
    $\sum_jK_j^\dagger K_j=\mathbb 1$. A \emph{channel layer} is a composition
    of channels on disjoint bonds, and a
    \emph{local channel circuit of depth $T$} is a composition
    $\Phi=\Phi_T\circ\dots\circ\Phi_1$ of $T$ channel layers.
    Its dual $\Phi^\dagger=\Phi_1^\dagger\circ\dots\circ\Phi_T^\dagger$
    replaces each channel on a bond by $A\mapsto\sum_jK_j^\dagger AK_j$. A
    density matrix is \emph{prepared in depth $T$ by local channels} if it
    equals $\Phi(\bigotimes_i\sigma_i)$ for a local channel circuit $\Phi$ of
    depth $T$ and density matrices $\sigma_i$ on single sites.
\end{definition}

\begin{lemma}[Duality and complete positivity of local channel circuits]
    \label{lem:qc_local_channel_circuit_properties}
    \uses{def:qc_local_channel_circuit, def:qc_local_circuit}
    \lean{QuantumCircuit.channelCircuitMap_isKrausCPTP,
        QuantumCircuit.trace_channelCircuitMap_mul,
        QuantumCircuit.channelCircuitMap_map_toChannelLayer}
    \leanok
    Let $\Phi$ be a local channel circuit. Then
    \begin{align}
        \Tr(\Phi(\rho)A) & =\Tr(\rho\,\Phi^\dagger(A))
        \label{eq:qc_channel_circuit_duality}
    \end{align}
    for all operators $\rho$ and $A$ on the sites, and $\Phi$ is trace
    preserving and completely positive. A local circuit $U$ of depth $T$ is
    the local channel circuit $\rho\mapsto U\rho U^\dagger$ of depth $T$ with
    one operator per channel.
\end{lemma}
\begin{proof}
    \leanok
    For one channel on a bond,
    $\Tr\bigl(\sum_jK_j\rho K_j^\dagger A\bigr)
        =\Tr\bigl(\rho\sum_jK_j^\dagger AK_j\bigr)$ by cyclicity of the
    trace, and $\sum_jK_j^\dagger K_j=\mathbb 1$ makes it trace preserving
    and completely positive. Both properties pass to compositions, hence to
    layers and circuits. A gate $G$ on a bond is the channel with the single
    operator $G$.
\end{proof}

\begin{lemma}[Light cone of a local channel circuit]
    \label{lem:qc_channel_light_cone}
    \uses{def:qc_local_channel_circuit, def:qc_light_cone_geometry,
        lem:qc_light_cone_geometries}
    \lean{QuantumCircuit.channelCircuitDual_mem_supportedOperators,
        QuantumCircuit.channelCircuitDual_mem_supportedOperators_lightCone,
        QuantumCircuit.BondChannelLayer.dual_mem_supportedOperators,
        QuantumCircuit.ChannelLayer.dual_mem_supportedOperators}
    \leanok
    Let $\Phi$ be a local channel circuit of depth $T$ and let $Z$ act on a
    set of sites $X$. Then $\Phi^\dagger(Z)$ acts on the light cone of radius
    $T$ of $X$; on the ring, on the sites within ring distance $T$ of $X$.
\end{lemma}
\begin{proof}
    \leanok
    By induction on the layers it suffices to treat one layer. The dual of a
    channel on a bond disjoint from $X$ fixes $Z$: its operators commute with
    $Z$, so $\sum_jK_j^\dagger ZK_j=\sum_jK_j^\dagger K_jZ=Z$. The dual of a
    channel on a bond meeting $X$ maps operators acting on the one-step
    neighbourhood of $X$ to such operators, since that set contains the bond.
    The channels of a layer commute, so the dual of the layer is the dual of
    the channels meeting $X$ applied after the dual of the others.
\end{proof}

\begin{lemma}[No connected correlations beyond the light cone of channels]
    \label{lem:qc_channel_light_cone_correlations}
    \uses{lem:qc_channel_light_cone, lem:qc_local_channel_circuit_properties}
    \lean{QuantumCircuit.trace_mul_mul_eq_of_isChannelPreparedInDepth,
        QuantumCircuit.IsBondChannelPreparedInDepth.trace_mul_mul_eq,
        QuantumCircuit.channelCircuitDual_mul,
        QuantumCircuit.channelCircuitDual_mul_of_disjoint_lightCone,
        QuantumCircuit.BondChannelLayer.dual_mul,
        QuantumCircuit.ChannelLayer.dual_mul,
        QuantumCircuit.trace_rectKronecker_mul_mul}
    \leanok
    Let $\rho$ be prepared in depth $T$ by local channels, and let $Z$, $Z'$
    act on sets of sites whose light cones of radius $T$ are disjoint, for
    instance on the ring sets at distance larger than $2T$. Then
    $\Tr(\rho ZZ')
        =\Tr(\rho Z)\Tr(\rho Z')$.
\end{lemma}
\begin{proof}
    \leanok
    The dual of a channel is not multiplicative, but the dual of one layer is
    multiplicative on operators $Z$, $Z'$ acting on sets $X$, $Y$ whose
    one-step neighbourhoods are disjoint. Indeed, no bond of the layer
    meets both $X$ and $Y$. If a channel acts off a set on which an operator
    $B$ acts, then its dual $\Psi$ satisfies $\Psi(AB)=\Psi(A)B$ and
    $\Psi(BA)=B\Psi(A)$. Writing the dual of the layer as the dual of the
    channels meeting $X$ after the dual of the others, and using
    Lemma~\ref{lem:qc_channel_light_cone} for the latter applied to $Z'$,
    one obtains the dual of $ZZ'$ as the product of the duals. By induction
    on the layers, $\Phi^\dagger(ZZ')=\Phi^\dagger(Z)\Phi^\dagger(Z')$, and
    the two factors act on disjoint sets. For a product density
    $\sigma=\bigotimes_i\sigma_i$ of trace one and products
    $\bigotimes_im_i$, $\bigotimes_im'_i$ with disjoint supports,
    $\Tr(\sigma\bigotimes_im_im'_i)
        =\prod_i\Tr(\sigma_im_im'_i)$, and at each site one of
    $m_i$, $m'_i$ is the identity; bilinearity and the duality
    $\Tr(\Phi(\sigma)A)=\Tr(\sigma\Phi^\dagger(A))$
    conclude.
\end{proof}

Ancillas are attached and discarded site by site, so the operations that
change the local dimension act on single sites.

\begin{definition}[Onsite channel]
    \label{def:qc_onsite_channel}
    \lean{Matrix.rectKronecker, Matrix.rectKronecker_one,
        QuantumCircuit.OnsiteChannel,
        QuantumCircuit.OnsiteChannel.map, QuantumCircuit.OnsiteChannel.dual,
        QuantumCircuit.OnsiteChannel.id, QuantumCircuit.OnsiteChannel.attach,
        QuantumCircuit.OnsiteChannel.attachZero,
        QuantumCircuit.OnsiteChannel.attach_siteMap_apply,
        QuantumCircuit.OnsiteChannel.discard,
        QuantumCircuit.OnsiteChannel.discard_siteMap_apply}
    \leanok
    An \emph{onsite channel} from sites of dimension $d$ to sites of dimension
    $e$ is given, at each site $i$, by operators $K_{ij}\colon\C^d\to\C^e$ with
    $\sum_jK_{ij}^\dagger K_{ij}=\mathbb 1$. It is the channel
    $\rho\mapsto\sum_JK_J\rho K_J^\dagger$ on any finite set of sites, such as
    the ring, with
    $K_J=\bigotimes_iK_{iJ(i)}$ for every choice $J$ of one index per site;
    its dual is $A\mapsto\sum_JK_J^\dagger AK_J$. Write
    $\Phi_i(\sigma)=\sum_jK_{ij}\sigma K_{ij}^\dagger$ for its action at site
    $i$. Attaching an ancilla of dimension $a$ in the state
    $\tau=\sum_kw_kw_k^\dagger$ with $\sum_k\|w_k\|^2=1$ is the onsite channel
    $X\mapsto X\otimes\tau$ from dimension $d$ to dimension $da$, and
    discarding it is the onsite channel given by the partial trace over the
    second factor of $\C^d\otimes\C^a$. Ancillas attached to each site and
    local operations on a site and its ancillas are the free operations of
    \cite{Piroli2021LOCC}.
\end{definition}

\begin{lemma}[Onsite channels are product channels]
    \label{lem:qc_onsite_channel_properties}
    \uses{def:qc_onsite_channel}
    \lean{QuantumCircuit.OnsiteChannel.map_isKrausCPTP,
        QuantumCircuit.OnsiteChannel.map_rectKronecker,
        QuantumCircuit.OnsiteChannel.map_finKronecker,
        QuantumCircuit.OnsiteChannel.trace_map_mul,
        QuantumCircuit.OnsiteChannel.discard_siteMap_attach_siteMap}
    \leanok
    An onsite channel $\Phi$ is trace preserving and completely positive, it
    maps $\bigotimes_i\sigma_i$ to $\bigotimes_i\Phi_i(\sigma_i)$, it satisfies
    $\Tr(\Phi(\rho)A)=\Tr(\rho\,\Phi^\dagger(A))$, and discarding an attached
    ancilla gives back the original operator.
\end{lemma}
\begin{proof}
    \leanok
    The operators $K_J$ satisfy
    $\sum_JK_J^\dagger K_J=\bigotimes_i\sum_jK_{ij}^\dagger K_{ij}=\mathbb 1$,
    so the Kraus form gives complete positivity and trace preservation. Since
    $K_J$ is a tensor product, $K_J(\bigotimes_i\sigma_i)K_J^\dagger
        =\bigotimes_iK_{iJ(i)}\sigma_iK_{iJ(i)}^\dagger$, and summing over $J$
    factorizes site by site. The duality is cyclicity of the trace. Finally,
    the partial trace of $X\otimes\tau$ is $\Tr(\tau)X=X$.
\end{proof}

\begin{lemma}[Onsite channels do not spread supports]
    \label{lem:qc_onsite_channel_locality}
    \uses{def:qc_onsite_channel}
    \lean{QuantumCircuit.OnsiteChannel.dual_mem_supportedOperators,
        QuantumCircuit.OnsiteChannel.dual_mul,
        QuantumCircuit.OnsiteChannel.dual_rectKronecker,
        QuantumCircuit.OnsiteChannel.dual_finKronecker,
        QuantumCircuit.rectangularKrausMap_rectKronecker_rectKronecker}
    \leanok
    Let $\Phi$ be an onsite channel. If $Z$ acts on a set of sites $X$, then
    $\Phi^\dagger(Z)$ acts on $X$. If moreover $Z'$ acts on a set of sites
    disjoint from $X$, then $\Phi^\dagger(ZZ')=\Phi^\dagger(Z)\Phi^\dagger(Z')$.
    The dual maps $\bigotimes_im_i$ to $\bigotimes_i\Phi_i^\dagger(m_i)$: for
    any operators $K_{ij}\colon\C^d\to\C^e$, the map
    $\rho\mapsto\sum_JK_J\rho K_J^\dagger$ with $K_J=\bigotimes_iK_{iJ(i)}$
    sends $\bigotimes_im_i$ to $\bigotimes_i\sum_jK_{ij}m_iK_{ij}^\dagger$.
\end{lemma}
\begin{proof}
    \leanok
    The dual maps $\bigotimes_im_i$ to $\bigotimes_i\Phi_i^\dagger(m_i)$, and
    $\Phi_i^\dagger(\mathbb 1)=\sum_jK_{ij}^\dagger K_{ij}=\mathbb 1$. Hence
    the identity factors off $X$ are kept. For products with disjoint
    supports, one of $m_i$, $m'_i$ is the identity at each site, so
    $\Phi_i^\dagger(m_im'_i)=\Phi_i^\dagger(m_i)\Phi_i^\dagger(m'_i)$.
    Linearity concludes.
\end{proof}

\begin{definition}[Local channel conversion]
    \label{def:qc_local_channel_conversion}
    \uses{def:qc_local_channel_circuit, def:qc_onsite_channel}
    \lean{QuantumCircuit.IsLocalChannelProtocol,
        QuantumCircuit.IsLocalChannelConversion}
    \leanok
    A \emph{local channel protocol of depth $T$} is a composition
    $\Psi=\Phi_T\circ L_T\circ\Phi_{T-1}\circ\dots\circ L_1\circ\Phi_0$ of
    channel layers $L_t$ and maps $\Phi_t$, each a composition of finitely
    many onsite channels, which may change the local dimension; $\Phi_0$
    contains at least one onsite channel, and $\Phi_t$ for $t\ge1$ may be
    the identity. A matrix $\rho$ on the chain is
    \emph{converted in depth $T$} into a matrix $\sigma$, possibly of
    another local dimension, if $\sigma=\Psi(\rho)$ for a local channel
    protocol $\Psi$ of depth at most $T$. This is the channel form of the circuits
    $V'=U_\ell V_\ell\cdots U_1V_1U_0$ of \cite{Piroli2021LOCC}, with
    ancillas attached to each site, local operations between the layers and
    the ancillas traced out at the end, without measurements.
\end{definition}

\begin{lemma}[Basic properties of local channel conversions]
    \label{lem:qc_local_channel_conversion_properties}
    \uses{def:qc_local_channel_conversion, def:qc_local_channel_circuit,
        lem:qc_onsite_channel_properties}
    \lean{QuantumCircuit.IsLocalChannelProtocol.isKrausCPTP,
        QuantumCircuit.IsLocalChannelProtocol.channelCircuitMap_comp,
        QuantumCircuit.isLocalChannelConversion_circuit,
        QuantumCircuit.IsChannelPreparedInDepth.exists_isLocalChannelConversion,
        QuantumCircuit.IsLocalChannelConversion.refl,
        QuantumCircuit.IsLocalChannelConversion.mono,
        QuantumCircuit.IsLocalChannelConversion.exists_isKrausCPTP,
        QuantumCircuit.IsLocalChannelConversion.density}
    \leanok
    A local channel protocol is trace preserving and completely positive.
    Attaching ancillas, running a local channel circuit of depth $T$ and
    discarding is a local channel protocol of depth $T$. Every density matrix
    is converted into itself in depth $0$, a conversion in depth $T$ is a
    conversion in every larger depth, a conversion of a density matrix is a
    density matrix, and a density matrix prepared in depth $T$ by local
    channels is converted in depth $T$ from a product density.
\end{lemma}
\begin{proof}
    \leanok
    By induction on the protocol: onsite channels are trace preserving and
    completely positive by Lemma~\ref{lem:qc_onsite_channel_properties},
    channel layers are by construction, and both properties are stable under
    composition. Appending the layers of a circuit of depth $T$ one at a time,
    each onsite channel keeps the depth and each layer raises it by one; the
    attaching and discarding channels are onsite. The identity onsite channel
    gives depth $0$, depth bounds are monotone, and a conversion is realized by
    a channel, which maps a positive semidefinite matrix of trace one to a
    positive semidefinite matrix of trace one. A preparation in depth $T$ by
    local channels starts from a product density and is itself such a
    protocol.
\end{proof}

\begin{lemma}[Composition of conversions]
    \label{lem:qc_local_channel_conversion_trans}
    \uses{def:qc_local_channel_conversion}
    \lean{QuantumCircuit.IsLocalChannelConversion.trans,
        QuantumCircuit.IsLocalChannelProtocol.comp}
    \leanok
    If $\rho$ is converted in depth $T_1$ into $\sigma$ and $\sigma$ is
    converted in depth $T_2$ into $\tau$, then $\rho$ is converted in depth
    $T_1+T_2$ into $\tau$.
\end{lemma}
\begin{proof}
    \leanok
    The composition $\Psi_2\circ\Psi_1$ of local channel protocols of depths
    $T_1$ and $T_2$ is a local channel protocol of depth $T_1+T_2$. Write
    $\Psi_2$ as the sequence of its onsite channels and channel layers, and
    apply them after $\Psi_1$ one at a time: the first onsite channel of
    $\Psi_2$ becomes a further onsite channel in the last group $\Phi_{T_1}$
    of $\Psi_1$, each onsite channel keeps the depth, and each of the $T_2$
    channel layers raises it by one.
\end{proof}

\begin{lemma}[Light cone of a local channel protocol]
    \label{lem:qc_protocol_light_cone}
    \uses{def:qc_local_channel_conversion, lem:qc_onsite_channel_locality,
        lem:qc_channel_light_cone, lem:qc_channel_light_cone_correlations}
    \lean{QuantumCircuit.IsLocalChannelProtocol.exists_dual}
    \leanok
    A local channel protocol $\Psi$ of depth $T$ has a dual $\Psi^\dagger$ with
    $\Tr(\Psi(\rho)A)=\Tr(\rho\,\Psi^\dagger(A))$ that maps operators acting
    on a set of sites $X$ to operators acting on the sites within distance
    $T$ of $X$, and that satisfies
    \begin{align}
        \Psi^\dagger(ZZ') & =\Psi^\dagger(Z)\,\Psi^\dagger(Z')
        \label{eq:qc_conversion_dual_mul}
    \end{align}
    whenever $Z$ and $Z'$ act on sets whose neighbourhoods of radius $T$ are
    disjoint. This is the light cone of a depth-$T$ circuit in the proof of
    \cite[Theorem~1]{Malz2024Preparation}, here with channels and onsite
    operations.
\end{lemma}
\begin{proof}
    \leanok
    By induction on the protocol. For an onsite channel this is
    Lemma~\ref{lem:qc_onsite_channel_locality} with radius $0$; a channel
    layer adds one to the radius, by Lemma~\ref{lem:qc_channel_light_cone}
    for the support and as in the proof of
    Lemma~\ref{lem:qc_channel_light_cone_correlations} for
    (\ref{eq:qc_conversion_dual_mul}), and an onsite channel keeps it.
\end{proof}

\begin{lemma}[No connected correlations after a conversion]
    \label{lem:qc_local_channel_conversion_correlations}
    \uses{def:qc_local_channel_conversion, lem:qc_protocol_light_cone}
    \lean{QuantumCircuit.trace_mul_mul_eq_of_isLocalChannelConversion}
    \leanok
    Let $\sigma$ be converted in depth $T$ from a product
    $\bigotimes_i\sigma_i$ of one-site matrices of trace one, and let $Z$,
    $Z'$ act on sets of sites at distance larger than $2T$. Then
    $\Tr(\sigma ZZ')=\Tr(\sigma Z)\Tr(\sigma Z')$. For pure states prepared
    from a product state by local unitary circuits this is the necessary
    condition for preparation by quantum circuits of \cite{Piroli2021LOCC}.
\end{lemma}
\begin{proof}
    \leanok
    Write $\rho=\bigotimes_i\sigma_i$ and $\sigma=\Psi(\rho)$ for a local
    channel protocol $\Psi$ of depth $T'\le T$, and let $\Psi^\dagger$ be its
    dual from Lemma~\ref{lem:qc_protocol_light_cone}, of radius $T'$.
    If $Z$ and $Z'$ act on sets at distance larger than $2T\ge2T'$, then
    $\Psi^\dagger(Z)$ and $\Psi^\dagger(Z')$ act on disjoint sets, and
    (\ref{eq:qc_conversion_dual_mul}) together with the factorization of the
    product density $\rho$ gives
    \begin{align}
        \Tr(\sigma ZZ') & =\Tr\bigl(\rho\,\Psi^\dagger(ZZ')\bigr)
        =\Tr\bigl(\rho\,\Psi^\dagger(Z)\Psi^\dagger(Z')\bigr) \notag \\
                        & =\Tr\bigl(\rho\,\Psi^\dagger(Z)\bigr)
        \Tr\bigl(\rho\,\Psi^\dagger(Z')\bigr)
        =\Tr(\sigma Z)\Tr(\sigma Z'). \notag
    \end{align}
\end{proof}

\section{Approximate local channel conversions}

The trace-norm approximation of~\cite[paragraph ``Phases of matter'']{Piroli2021LOCC}
can be applied to the local channels of
Definition~\ref{def:qc_local_channel_conversion}. This restricts the allowed
operations: the source additionally permits measurements and classical
feedforward, whereas the channels in this section do not; see
\cite{gap:psc21_local_channel_phase_scope}.

\begin{definition}[Approximate local channel conversion]
    \label{def:qc_approx_local_channel_conversion}
    \uses{def:qc_local_channel_conversion}
    \lean{QuantumCircuit.chainTraceNorm,
        QuantumCircuit.IsApproxLocalChannelConversion}
    \leanok
    A chain operator $\rho$ is approximately converted into $\sigma$ in
    depth $T$ with error $\varepsilon$ if there is a local channel protocol
    $\Psi$ of depth at most $T$ with
    $\lVert\Psi(\rho)-\sigma\rVert_1\le\varepsilon$.
    Here $\lVert A\rVert_1$ is the sum of the singular values of $A$.
\end{definition}

\begin{theorem}[Zero error and monotonicity]
    \label{thm:qc_approx_conversion_properties}
    \uses{def:qc_approx_local_channel_conversion}
    \lean{QuantumCircuit.IsLocalChannelConversion.approx,
        QuantumCircuit.IsApproxLocalChannelConversion.refl,
        QuantumCircuit.IsApproxLocalChannelConversion.mono,
        QuantumCircuit.isApproxLocalChannelConversion_zero_iff}
    \leanok
    Every operator is converted into itself with depth and error zero.
    Increasing the depth or error bound preserves approximate conversion.
    Approximate conversion with error zero is equivalent to exact conversion
    with the same depth bound.
\end{theorem}
\begin{proof}
    \leanok
    Use the identity channel, enlarge the bounds, and use
    $\lVert A\rVert_1=0$ if and only if $A=0$, respectively.
\end{proof}

\begin{theorem}[Composition of approximate conversions]
    \label{thm:qc_approx_conversion_trans}
    \uses{def:qc_approx_local_channel_conversion,
        lem:qc_local_channel_conversion_trans}
    \lean{QuantumCircuit.IsLocalChannelProtocol.chainTraceNorm_map_le,
        QuantumCircuit.IsApproxLocalChannelConversion.trans}
    \leanok
    Let $\rho$ and $\sigma$ be Hermitian. If $\rho$ is approximately converted
    into $\sigma$ with depth $T_1$ and error $\varepsilon_1$, and $\sigma$ into
    $\tau$ with depth $T_2$ and error $\varepsilon_2$, then $\rho$ is
    approximately converted into $\tau$ with depth $T_1+T_2$ and error
    $\varepsilon_1+\varepsilon_2$.
\end{theorem}
\begin{proof}
    \leanok
    The composition $\Psi_2\circ\Psi_1$ has the sum of the two depths.
    The error $\Psi_1(\rho)-\sigma$ is Hermitian. Trace-norm contractivity of
    trace-preserving positive maps and the triangle inequality give
    \begin{align}
      \lVert\Psi_2(\Psi_1(\rho))-\tau\rVert_1
      &\le \lVert\Psi_2(\Psi_1(\rho)-\sigma)\rVert_1
          +\lVert\Psi_2(\sigma)-\tau\rVert_1 \notag\\
      &\le \lVert\Psi_1(\rho)-\sigma\rVert_1
          +\lVert\Psi_2(\sigma)-\tau\rVert_1
      \le\varepsilon_1+\varepsilon_2. \notag
    \end{align}
\end{proof}

\begin{definition}[Asymptotic local channel conversion]
    \label{def:qc_asymptotic_local_channel_conversion}
    \uses{def:qc_approx_local_channel_conversion}
    \lean{QuantumCircuit.IsAsymptoticLocalChannelConversion,
        QuantumCircuit.IsLocalChannelPhaseEquivalent}
    \leanok
    Let $\rho_N$ and $\sigma_N$ be families of operators on chains of $N$
    sites, with fixed local dimensions that may differ between the families.
    The family $\rho_N$ is asymptotically locally convertible into $\sigma_N$
    if there are an integer $p\ge0$, a constant $C\ge0$, and errors
    $\varepsilon_N\to0$ such that, for all sufficiently large $N$, there
    is an approximate conversion with error $\varepsilon_N$ and depth
    $T_N\le C(1+\log(N+1))^p$. Two families are locally phase equivalent
    if they are asymptotically locally convertible in both directions.
\end{definition}

\begin{theorem}[Local channel phase equivalence]
    \label{thm:qc_local_channel_phase_equivalence}
    \uses{def:qc_asymptotic_local_channel_conversion,
        thm:qc_approx_conversion_trans}
    \lean{QuantumCircuit.IsAsymptoticLocalChannelConversion.refl,
        QuantumCircuit.IsAsymptoticLocalChannelConversion.trans,
        QuantumCircuit.IsLocalChannelPhaseEquivalent.refl,
        QuantumCircuit.IsLocalChannelPhaseEquivalent.symm,
        QuantumCircuit.IsLocalChannelPhaseEquivalent.trans,
        QuantumCircuit.IsLocalChannelPhaseEquivalent.equivalence}
    \leanok
    Local channel phase equivalence is an equivalence relation on eventually
    Hermitian families, and in particular on density-matrix families.
\end{theorem}
\begin{proof}
    \leanok
    Identity channels give reflexivity; the definition gives symmetry.
    For transitivity, compose the approximate conversions in each direction.
    The sum of the errors tends to zero. For $b_N=1+\log(N+1)\ge1$, the
    depth satisfies
    \begin{align}
      C_1b_N^{p_1}+C_2b_N^{p_2}
      &\le(C_1+C_2)b_N^{\max(p_1,p_2)}. \notag
    \end{align}
    Thus the depth remains polylogarithmic.
\end{proof}

